\setcounter{section}{3}
\setcounter{theorem}{12}
\section{A Natural Complete Problem}
\label{sec:chevalley}





We start with some notation that will be useful for the
presentation of our results.
\smallskip

\paragr{Notations.} For any polynomial $g \in \bbF_p[\vecx]$, we
define $\deg(g)$ to be the degree of $g$. We define
\textit{the expansion to monic monomials} of $g$ as $\sum_{\ell = 1}^L t_{\ell}(\vecx)$, 
where $t_{\ell}(\vecx)$ is a monic monomial in
$\bbF_p[\vecx]$, i.e. a monomial with coefficient $1$. For example,
the expansion of the polynomial
$g(x_1, x_2) = x_1 \cdot (2 x_1 + 3 x_2)$ is
given by $x_1^2 + x_1^2 + x_1 x_2 + x_1 x_2 + x_1 x_2$.

For a polynomial system $\vec{f} := (f_1, \ldots, f_m) \in \bbF_p[\vecx]^m$, its affine variety $\calV_{\vecf} \subseteq \bbF_p^n$ is defined as
$\calV_{\vecf} \coloneqq \set{\vecx \in \bbF_p^n \mid \vecf(\vecx) = \vec{0}}$.
Let $\overline{\calV_{\vecf}} \coloneqq \bbF_p^n \setminus \calV_{\vecf}$. If the constant term of each $f_i$ is $0$, we say that $\vecf$ is {\em zecote}, standing for ``Zero Constant Term'' (owing to lack of known terminology and creativity on our part).

  
\subsection{The Chevalley-Warning Theorem} \label{sec:ChevalleyWarningProof}

  We repeat the formal statement of Chevalley-Warning Theorem together with its
proof.

\begin{cwthm}[\cite{r15,r37}]
    For any prime $p$ and a polynomial system
  $\vecf \in \bbF_p[\vecx]^m$ satisfying $\sum_{i = 1}^m \deg(f_i) < n$
  \emph{(CW Condition)}, $|\calV_{\vecf}| \equiv 0 \mod{p}$.
\end{cwthm}

  We describe the proof of \CWT through \autoref{lem:mainLemmaInProofOfCWT}.
Even though there are direct proofs, the following presentation helps motivate the generalizations we study in future sections. Given a polynomial system $\vecf \in
\bbF_p[\vecx]^m$, a key idea in the proof is the polynomial $\CW_{\vecf}(\vecx)
\coloneqq \prod_{i = 1}^m \CW_{f_i}(\vecx)$ where each $\CW_{f_i}(\vecx) := (1 - f_i(\vecx)^{p - 1})$. Observe that $\CW_{\vecf}(\vecx) = 1$ if
$\vecx \in \calV_{\vecf}$ and is $0$ otherwise. The following definition
describes the notion of a max-degree monomial of $\CW_{\vecf}$ that plays an important role in the proof.

\begin{definition}[\textsc{Max-Degree Monic Monomials}]
\label{def:max-degreeMonomial}
For any prime $p$, let $\vecf \in \bbF_p[\vecx]^m$ and let the expansion into monic monomials of $\CW_{f_i}(\vecx)$ be $\sum_{\ell = 1}^{r_i} t_{i, \ell}(\vecx)$. Let also $U_i = \{(i, \ell) \mid \ell \in [r_i]\}$ and $U = \bigtimes_{i = 1}^m U_i$, we define the following quantities.
  \begin{Enumerate}
    \item A \textit{monic monomial} of $\CW_{\vecf}$ is a product $t_S(\vecx) = \prod_{i = 1}^m t_{s_i}(\vecx)$ for $S = (s_1, \ldots, s_m) \in U$.
	    \item A \textit{max-degree monic monomial} of $\CW_{\vecf}$ is any monic monomial $t_S(\vecx)$, such that\\
    $t_S(\vecx) \equiv \prod_{j = 1}^n x_j^{p - 1} \mod{\set{x_i^p - x_i}_{i \in [n]}}$.
    \item We define $\calM_{\vecf}$ to be the set of max-degree monic monomials of $\CW_{\vecf}$, i.e. \\
	$\calM_{\vecf} \coloneqq \set{S \in U \mid t_S \text{ is a max-degree monic monomial of } \CW_{\vecf}}$.
  \end{Enumerate}
\end{definition}

\noindent In words, the monomials $t(S)$ are precisely the ones that arise when symbolically expanding $\CW_{\vecf}(\bx)$. We illustrate this with an example: Let $p = 3$ and $f_1(x_1, x_2) = x_1 + x_2$ and $f_2(x_1, x_2) = x_1^2$. Then modulo $\set{x_1^3 - x_1, x_2^3 - x_2}$, we have
\begin{align*}
\CW_{(f_1,f_2)}(x_1, x_2) &~=~ (1 - (x_1 + x_2)^2)(1 - (x_1^2)^2)\\
&~=~ (1 - x_1^2 - 2 x_1 x_2 - x_2^2) \cdot (1 - x_1^2)\\
&~=~ (1 + x_1^2 + x_1^2 + x_1x_2 +x_2^2 + x_2^2) \cdot (1 + x_1^2 + x_1^2)
\end{align*}

\noindent Thus there are $18$ $(= 6 \times 3)$ monic monomials in the system $(f_1, f_2)$. The monomial corresponding to $S = ((1,5),(2,2))$ is a maximal monomial since the $5$-th term in $\CW_{f_1}$ is $x_2^2$ and $2$-nd term in $\CW_{f_2}$ is $x_1^2$. Using the above definitions, we now state the main technical lemma of the proof of \CWT.

\begin{lemma}[Main Lemma in the proof of CWT] \label{lem:mainLemmaInProofOfCWT}
   For any prime $p$ and any system of polynomials
  $\vecf \in \bbF_p[\vecx]^m$, it holds that
  $\abs{\calV_{\vecf}} \equiv (-1)^n \abs{\calM_{\vecf}} \mod{p}$.
\end{lemma}

\begin{proof}
As noted earlier, $\CW_{\vecf}(\vecx) = 1$ if $\bx \in \calV_{\vecf}$ and is $0$ otherwise. Thus, it follows that $\abs{\calV_{\vecf}} \equiv \sum_{\bx \in \bbF_p^n} \CW_{\vecf}(\bx) \mod{p}$. For any monic monomial $m(\bx) = \prod_{j=1}^n x_j^{d_j}$, it holds that $\sum_{\bx \in \bbF_p^n} m(\bx) = 0$ if $d_j < p - 1$ for some $x_j$. On the other hand, for the monic max-degree monomial $m(\bx) = \prod_{j = 1}^n x_j^{p-1}$, it holds that $\sum_{\bx \in \bbF_p^n} m(\bx) = (p-1)^n$. Thus, we get that $|\calV_{\vecf}| \equiv \sum_{\bx \in \bbF_p^n} \CW_{\vecf}(\bx) \mod{p} \equiv \sum_{S\in U} \sum_{\bx \in \bbF_p^n} t_S(\vecx)  \mod{p} \equiv (-1)^n |\calM_{\vecf}| \mod{p}$.
\end{proof}

\noindent The proof of Chevalley-Warning Theorem follows easily from
\autoref{lem:mainLemmaInProofOfCWT}.

\begin{proof}[Proof of Chevalley-Warning Theorem.]
We have that $\deg(\CW_{\vecf}) \le (p - 1) \sum_{i = 1}^m \deg(f_i)$. Thus, if $\vecf$ satisfies \eqref{eq:CWcondition}, then $\deg(\CW_{\vecf}) < (p - 1) n$ and hence $\abs{\calM_{\vecf}} = 0$. \CWT now follows from \autoref{lem:mainLemmaInProofOfCWT}.
\end{proof}

\subsection{The Chevalley-Warning Theorem with Symmetry}
\label{sec:ChevalleyWarningWithSymmetry}

  In this section, we formalize the intuition that we built in
\hyperref[sec:proofsOfCancellation]{Sections~\ref{sec:proofsOfCancellation}} and \ref{sec:computationalChevalleyWarning} to prove the more general
statements to lead to the same conclusion as the Chevalley-Warning
Theorem.

  First, we prove a theorem that argues about the cardinality of $\calV_{\vecf}$
directly using some symmetry of the system of polynomials $\vecf$. Then,
combining this symmetry-based argument with the \eqref{eq:generalCWCondition} we get
the generalization of the Chevalley-Warning Theorem. Our natural
$\PPA_p$-complete problem is based on this generalization.

The theorem statements are simplified using the definition of
\textit{free action} of a group. For a permutation over $n$ elements $\sigma \in S_n$, we define
$\langle \sigma \rangle$ to be the sub-group generated by $\sigma$ and
$|\sigma|$ to be the order of $\langle\sigma \rangle$. For $\vecx \in \bbF_p^n$, $\sigma(\vecx)$ denotes the assignment obtained by permutating the variables of the assignment $\vecx$ according to $\sigma$.

\begin{definition}[\textsc{Free Group Action}]
    Let $\sigma \in S_n$ and $\calV \subseteq \bbF_p^n$, then we say that
  $\langle \sigma \rangle$ acts freely on $\calV$ if, for every
  $\vecx \in \calV$, it holds that $\sigma(\vecx) \in \calV$ and
  $\vecx \neq \sigma(\vecx)$.
\end{definition}

  Our first theorem highlights the
use of symmetry in arguing about the size of $\abs{\calV_{\vecf}}$.

\begin{theorem} \label{thm:symmetricChevalleyTheorem}
    Let $\vecf \in \F_p[\vecx]^m$ be a system of polynomials.
  If there exists a permutation $\sigma \in S_n$ with
  $\abs{\sigma} = p$ such that $\a{\sigma}$ acts freely on
  $\overline{\calV}_{\vecf}$, then $\abs{\calV_{\vecf}} \equiv 0 \mod{p}$.
\end{theorem}

\begin{proof}
    Since $\sigma$ acts freely on $\overline{\calV}_{\vecf}$, we can partition
  $\overline{\calV}_{\vecf}$ into orbits of any
  $\vecx \in \overline{\calV}_{\vecf}$ under actions of $\a{\sigma}$, namely
  sets of the type $\set{\sigma^i(\vecx)}_{i \in [p]}$ for
  $\vecx \in \overline{\calV}_{\vecf}$. Since $\a{\sigma}$ acts freely on
  $\overline{\calV}_{\vecf}$, each such orbit has size $p$. Thus, we can
  conclude that $\abs{\overline{\calV}_{\vecf}} \equiv 0 \pmod p$ from which
  the theorem follows.
\end{proof}

\begin{remark} \label{rem:syntacticVerificationOfSymmetry}
  For any polynomial system $\vecf$ and any
  permutation $\sigma$, we can check in linear time if $\abs{\sigma} = p$ and we
  can syntactically refute that $\langle \sigma \rangle$ acts freely on
  $\overline{\calV}_{\vecf}$ with an $\vecx \in \bbF_p^n \smallsetminus \set{\vec{0}}$ such that $\vecf(\sigma(\vecx)) = \vec{0}$ or $\sigma(\vecx) = \vecx$.
\end{remark}
\bigskip

We now state and prove an extension of \CWT that captures both
the argument from \autoref{lem:mainLemmaInProofOfCWT} and the symmetry
argument from \autoref{thm:symmetricChevalleyTheorem}.

\begin{theorem}[\textsc{Chevalley-Warning with Symmetry Theorem}]
\label{thm:combinedChevalleyTheorem}
	  Let $\vecg \in \F_p[\vecx]^{m_g}$ and
	$\vech \in \F_p[\vecx]^{m_h}$ be two systems of polynomials, and
	$\vecf \coloneqq (\vecg, \vech)$. If there exists a permutation $\sigma \in S_n$ with $\abs{\sigma} = p$ such that (1) $\calM_{\vecg} = \emptyset$ and (2) $\a{\sigma}$ acts freely on
	$\calV_{\vecg} \cap \overline{\calV_{\vech}}$, then $\abs{\calV_{\vecf}} = 0 \mod{p}$.
\end{theorem}

\begin{remark} \label{rem:conditionsOfChevalleyWithSymmetry}
    We point to the special form of Condition 2. By definition,
  $\calV_{\vecf} = \calV_{\vecg} \cap \calV_{\vech}$, hence if $\a{\sigma}$
  were to act freely on $\overline{\calV_{\vecg}} \cup \overline{\calV_{\vech}}$ (or even $\calV_{\vecg} \cap \calV_{\vech}$), then we could just
  use \autoref{thm:symmetricChevalleyTheorem} to get that
  $\abs{\calV_{\vecf}} \equiv 0 \mod{p}$. In the above theorem,
  we only require that $\a{\sigma}$ acts freely on
  $\calV_{\vecg} \cap \overline{\calV_{\vech}}$. Observe that
  \autoref{thm:symmetricChevalleyTheorem} follows as a special case of \CWT with
  Symmetry by setting $m_g = 0$. Additionally, by setting $m_h = 0$ we get the generalization of \CWT corresponding to the
  \eqref{eq:generalCWCondition} as presented in
  \autoref{sec:computationalChevalleyWarning}.
\end{remark}

\begin{proof}[Proof of \autoref{thm:combinedChevalleyTheorem}]
  	If $\CW_{\vecg}$ does not have any max-
	degree monic monomials, we have $\abs{\calV_{\vecg}} \equiv 0 \mod{p}$ (similar to proof of \CWT) and, since
	$\overline{\calV_{\vecg}} = \F_p^n \setminus \calV_{\vecg}$, we have $\abs{\overline{\calV_{\vecg}}} \equiv 0 \mod{p}$. Also,
	since $\a{\sigma}$ acts freely on $\calV_{\vecg} \cap
	\overline{\calV_{\vech}}$, we have
	$\abs{\calV_{\vecg} \cap \overline{\calV_{\vech}}} \equiv 0 \mod{p}$ (similar to the proof of
	\autoref{thm:symmetricChevalleyTheorem}). Hence,
	$\abs{\overline{\calV_{\vecf}}} = \abs{\overline{\calV_{\vecg} \cap
	\calV_{\vech}}}  = \abs{\overline{\calV_{\vecg}} \cup
	\overline{\calV_{\vech}}} = \abs{\overline{\calV_{\vecg}}} +
	\abs{\calV_{\vecg} \cap \overline{\calV_{\vech}}} \equiv  0 \mod{p}$. Thus,
	$\abs{\calV_{\vecf}} \equiv 0 \mod{p}$.
\end{proof}

\subsection{Computational Problems Related to Chevalley-Warning Theorem}
\label{sec:computationalChevalleyWarningFormal}

  We now follow the intuition developed in the previous section and in
\autoref{sec:intro-complete} to formally define the computational problems
$\eChevalleyp$, $\lChevalleyp$, and $\cChevalleyp$.
\medskip


	\noindent
	\begin{minipage}{\linewidth}
		\begin{definition}
			\label{def:explicit-chevalley} \contour{black}{($\eChevalleyp$)}\vspace{-5pt}
			\begin{itemize}[leftmargin=2.15cm,itemsep=4pt]
				\item[\slshape Principle:] Chevalley-Warning Theorem.
				\item[\slshape Input:] $\vecf \in \bbF_p[\vecx]^m$ : an explicit zecote polynomial system.
				\item[\slshape Condition:] $\sum_{i = 1}^m \deg(f_i) < n$.
				\item[\slshape Output:] $\vecx \in \bbF_p^n$ such that $\vecx \neq \vec{0}$ and $\vecf(\vecx) = \vec{0}$.
			\end{itemize}
		\end{definition}
	\end{minipage}\vspace{3mm}

\bigskip


	\noindent
	\begin{minipage}{\linewidth}
		\begin{definition}
			\label{def:general-chevalley} \contour{black}{($\lChevalleyp$)}\vspace{-5pt}
			\begin{itemize}[leftmargin=2.15cm,itemsep=4pt]
				\item[\slshape Principle:] General Chevalley-Warning Theorem via \eqref{eq:generalCWCondition}.
				\item[\slshape Input:] $\vecf \in \bbF_p[\vecx]^m$ : an explicit zecote polynomial system.
				\item[\slshape Output:] \begin{Enumerate}[leftmargin=*]
  \item[0.] A max-degree monic monomial $t_S(\vecx)$ of $\CW_{\vecf}$, or
  \item $\vecx \in \bbF_p^n$ such that $\vecx \neq \vec{0}$ and $\vecf(\vecx) = \vec{0}$.
\end{Enumerate}
			\end{itemize}
		\end{definition}
	\end{minipage}\vspace{3mm}


\bigskip


	\noindent
	\begin{minipage}{\linewidth}
		\begin{definition}
			\label{def:combined-chevalley} \contour{black}{($\cChevalleyp$)}\vspace{-5pt}
			\begin{itemize}[leftmargin=2.15cm,itemsep=4pt]
				\item[\slshape Principle:] Chevalley-Warning Theorem with Symmetry (\autoref{thm:combinedChevalleyTheorem}).
				\item[\slshape Input:] $\triangleright$ $\vecg \in \bbF_p[\vecx]^{m_g}$ and $\vech \in \bbF_p[\vecx]^{m_h}$ : explicit zecote polynomial systems\\
 $\triangleright$ $\sigma \in S_n$ : a permutation over $[n]$.
				\item[\slshape Condition:] $|\sigma| = p$.
				\item[\slshape Output:] \begin{Enumerate}[leftmargin=*]
  \item[0.] (a) A max-degree monic monomial $t_S(\vecx)$ of $\CW_{\vecg}$, or
  \item[] (b) $\vecx \in \calV_{\vecg} \cap \overline{\calV_{\vech}}$ such that
        $\sigma(\vecx) \not\in (\calV_{\vecg} \cap \overline{\calV_{\vech}}) \smallsetminus \set{\vecx}$, or
  \item $\vecx \in \bbF_p^n$ such that $\vecx \neq \vec{0}$ and $\vecf(\vecx) = \vec{0}$.
\end{Enumerate}
			\end{itemize}
		\end{definition}
	\end{minipage}\vspace{3mm}


\begin{remark} \label{rem:computationalDefinitions}
    Some observations about the above computational problems follow:
  \begin{Enumerate}
    \item In the problems $\lChevalleyp$ and $\cChevalleyp$,  we assume that, if the output is a max-degree monic monomial, this is given via the multiset of indices $S$ that describes the monomial as formalized in \autoref{def:max-degreeMonomial}.
    \item We have $\eChevalleyp \preceq \lChevalleyp \preceq \cChevalleyp$. Thus, inclusion of $\cChevalleyp$ in $\PPA_p$ implies that the problems $\eChevalleyp$ and $\lChevalleyp$ are in $\PPA_p$. Also, in Section6 we prove that $\sis_p \preceq \eChevalleyp$, where $\sis_p$ is a cryptographically relevant problem. This shows that the $\lChevalleyp$ and the $\cChevalleyp$ problems are at least as hard as the $\sis_p$ problem.
  \end{Enumerate}
\end{remark}

\noindent We restate our main result.

\completeness*

\subsection{\boldmath $\cChevalleyp$ is $\PPAp$--complete}
\label{sec:proofOfCombinedChevalley}

  We fist prove that $\cChevalley$ is in $\PPAp$ and then
prove its $\PPAp$-hardness.


\subsubsection{\boldmath $\cChevalleyp$ is in $\PPA_p$}

Even though Papadimitriou~\cite{r33} provided a rough proof sketch of $\Chevalley_p \in \PPAp$, a formal proof was not given. We show that $\cChevalleyp$ is in $\PPAp$ (and so are $\lChevalleyp$ and $\eChevalleyp$). In order to do so we extend the definition of $\bipartiteq$ to instances where the vertices might have exponential degree and edges appear with multiplicity. The key here is to define a $\bipartite_q$ instance with unbounded (even exponential) degree, but with additional information that allows us to verify solutions efficiently.\\[-1mm]


\noindent
\begin{minipage}{\linewidth}
\begin{definition}
	\label{def:succ-bipartite} \contour{black}{($\sucBipartite_q$)}\vspace{-5pt}
	\begin{itemize}[leftmargin=2.15cm,itemsep=4pt]
		\item[\slshape Principle:] Similar to $\bipartite_q$, but degrees are allowed to be exponentially large, edges are allowed with multiplicities at most $q-1$.
		\item[\slshape Object:] Bipartite graph $G = (V \cup U, E)$ s.t. $E \subseteq V \times U \times \bbZ_q$. Designated edge $e^* \in E$.
		\item[\slshape Inputs:] Let $V \coloneqq \set{0} \times \bit^{n-1}$ and $U \coloneqq \set{1} \times \bit^{n-1}$:\\
	$\triangleright$ $\calC:V \times U \to [q]$, edge counting circuit\\
	$\triangleright$ $\phi_V: V \times U \times [q] \to (U \times [q])^{q}$, grouping pivoted at $V$\\
	$\triangleright$ $\phi_U: V \times U \times [q] \to (V \times [q])^{q}$, grouping pivoted at $U$\\
	$\triangleright$ $e^* = (v^*, u^*, k^*)$, designated edge
		\item[\slshape Encoding:] $V \coloneqq \set{0} \times \bit^{n-1}$, $U \coloneqq \set{1} \times \bit^{n-1}$,\\
	$E \coloneqq \set{(v,u,k) : 1 \le k \le C(v,u),\ (v,u) \in V \times U}$ (here $k$ distinguishes multiplicities)\\
	Edge $(v,u,k)$ is grouped with $\set{(v, u', k') : (u',k') \in \phi_V(v,u,k)}$ (pivoting at $v$),\\
	\hspace*{1cm} provided $|\phi_V(v,u,k)| = q$, all $(v,u',k') \in E$ and $\phi_V(v,u',k') = \phi_V(v,u,k)$.\\
	Edge $(v,u,k)$ is grouped with $\set{(v', u, k') : (v',k') \in \phi_U(v,u,k)}$ (pivoting at $u$),\\
	\hspace*{1cm} provided $|\phi_U(v,u,k)| = q$, all $(v,u',k') \in E$ and $\phi_U(v',u,k') = \phi_V(v,u,k)$.
		\item[\slshape Solutions:] $e^*$ if $e^*$ is grouped, pivoting at $v^*$, or if $e^*$ is not grouped pivoting at $u^*$, OR\\
	$e \ne e^*$ if $e$ is not grouped pivoting at one of its ends.
	\end{itemize}
\end{definition}
\end{minipage}\vspace{5mm}

In words, $\sucBipartite_p$ encodes a bipartite graph with arbitrary degree. Instead of listing the neighbors of a vertex using a circuit, we have a circuit that outputs the multiplicity of edges between any two given vertices. We are therefore unable to efficiently count the number of edges incident on any vertex. The grouping function $\phi_V$ aims to group edges incident on any vertex $v \in V$ into groups of size $q$. Similarly, $\phi_U$ aims to group edges incident on any vertex $u \in U$. The underlying principle is that if we have an edge $e^*$ that is not grouped pivoting at $v^*$ (one of its endpoints), then either $e^*$ is not pivoted at $u^*$ (its other endpoint) or there exists another edge that is also not grouped pivoting at one of its ends. Note that in contrast to the problems previously defined, $v^*$ might still be an endpoint of a valid solution.

\begin{lemma} \label{lem:combinedChevalleyPPAp}
 For all primes $p$, $\cChevalleyp \in \PPAp$.
\end{lemma}
\begin{proof}We reduce $\cChevalleyp$ to $\sucBipartite_p$, which we show to be $\PPA_p$--complete in \autoref{sec:sucBipartite}. Given an instance of $\cChevalleyp$, namely a zecote polynomial system $\vecf = (\vecg, \vech)$ and a permutation $\sigma$, we construct a bipartite graph $G = (U \cup V, E)$ encoded as an instance of $\sucBipartite_p$ as follows. 

 \paragraph{Description of vertices.} $U = \bbF_p^n$, namely all possible assignments
 of $\vecx$. The vertices of $V$ are divided into two parts $V_1 \cup V_2$.
 The part $V_1$ contains one vertex for each monomial in the expansion of $\CW_{\vecg} = \prod_{i = 1}^{m_g} (1 - g_i^{p-1})$. Since $p$ is constant, we can
 efficiently list out the monomials of $1 - g_{i}^{p-1}$. For a fixed lexicographic ordering of the
 monomials of each $\CW_{g_i} \coloneqq 1 - g_{i}^{p-1}$, a
 monomial of $\CW_{\vecg}$ is represented by a tuple
 $(a_1, a_{2}, \dots, a_{m_g})$ with $0 \leq a_{i} < L_{i}$, where
 $a_{i}$ represents the index of a monomial of $\CW_{g_i}$ and  $L_{i}$
 is the number of monomials of $\CW_{g_i}$, where $a_i = 0$ corresponds to the constant term $1$. The part $V_2 \coloneqq \binom{\bbF_p^n}{p}$, i.e. it contains a vertex for each subset of $p$ distinct elements in $\bbF_p^n$.


 \paragraph{Description of edges.} We first describe the edges between $U$ and $V_1$, namely include an
 edge between an assignment $\vecx$ and a monomial $t$ with
 multiplicity $t(\vecx)$. With these edges in place, the degree of vertices are as follows:\vspace{-2mm}
 \begin{itemize}[leftmargin=0.5cm]
	\item $\vecx = 0^n$ has a single edge corresponding to the constant monomial $1$, since $\vecf$ is zecote. We let this be the designated edge $e^*$ in the final $\sucBipartite_p$ instance.
	\item $\vecx \notin \calV_{\vecg}$ has $0 \pmod p$ edges
	(counting multiplicities). Since $\CW_{\vecg}(\vecx) = 0$, the sum over all monomials of $t(\vecx)$ must be $0 \mod{p}$.
	\item $\vecx \in \calV_{\vecg}$ has $1 \pmod p$ edges (counting multiplicities), since the sum over all $t(\vecx)$ monomials gives $\CW_{\vecg}(\bx) \equiv 1 \mod{p}$.
 \end{itemize}
 Thus with the edges so far, the vertices (excluding $0^n$), with degree $\not\equiv 0 \mod{p}$ are precisely vertices  $t \in V_1$ such that $\sum_{\vecx} t(\vecx) \not\equiv 0 \pmod p$ or $\vecx \in \calV_{\vecg} \smallsetminus \set{0^n}$. For the former case, if $t$ contained a variable with degree less than $p-1$, then $\sum_{\vecx} t(\vecx) \equiv 0 \pmod p$. Hence, it must be that $t = \prod\limits_{i = 1}^{ n} x_i^{p-1}$. In the later case, the degree of $\vecx $ is $1 \mod{p}$ and hence $\vecx \in \calV_{\vecg}$. 
 However, there is no guarantee that a vertex $\vecx$ with degree $1 \mod{p}$ is in $\calV_{\vech}$ as well. To argue
 about $\vech$, we add edges between $U$ and $V_2$ that exclude solutions $\vecx \in \calV_{\vecg} \cap \overline{\calV_{\vech}}$, on which $\sigma$ acts freely (that is, $\sigma(\vecx) = \vecx$). More specifically, for
 $\vecx \in \calV_{\vecg} \cap \overline{\calV_{\vech}}$, if $\sigma(\vecx) \ne \vecx$,
  we add an edge with multiplicity $p-1$ between $\vecx$ and $\Sigma_{\vecx} \in V_2$ where $\Sigma_{\vecx} \coloneqq \{\sigma^i(\vecx)\}_{i \in \bbZ_p}$ (note that, in this case $\abs{\Sigma_{\vecx}} = p$ since $\sigma(\vecx) \ne \vecx$ and $|\sigma| = p$ is prime).
 Observe that, if a vertex in $V_2$ corresponds
 to a $\Sigma_{\vecx}$, it has $p$ edges each with multiplicity $p-1$, one for each
 $\vecx' \in \Sigma_{\vecx}$ only if $\Sigma_{\vecx} \subseteq \calV_{\vecg} \cap \overline{\calV_{\vech}}$. If a vertex in $V_2$ does not correspond to a
 $\Sigma_{\vecx}$, then it has no edges. Thus, a vertex in $V_2$ has degree $\not\equiv 0 \mod{p}$ iff it contains an $\vecx \in \calV_{\vecg} \cap \overline{\calV_{\vech}}$ such that $\sigma(\vecx) \notin \calV_{\vecg} \cap \overline{\calV_{\vech}}$.

 Thus, with all the edges added, vertices with degree $\not\equiv 0 \mod{p}$ correspond to one of\vspace{-2mm}
 \begin{itemize}
 \item $\vecx \in \calV_{\vecg} \cap \calV_{\vech}$ such that $\vecx \ne \vec{0}$, or
 \item $t \in V_1$ such that $t(\vecx)$ is a max-degree monomial or
 \item $\vecx \in \calV_{\vecg} \cap \overline{\calV_{\vech}}$ such that $\sigma(\vecx) = \vecx$ or
 \item $v \in V_2$ such that $\exists\, \vecx \in v$ satisfying $\vecx \in \calV_{\vecg} \cap \overline{\calV_{\vech}}$ and $\sigma(\vecx) \notin \calV_{\vecg} \cap \overline{\calV_{\vech}}$.
 \end{itemize}
 These correspond precisely to the solutions of $\cChevalleyp$. To summarize, the edge
 counting circuit $C$ on input $(\vecx, t) \in U \times V_1$ outputs
 $t(\vecx)$ and on input $(\vecx, v) \in U \times V_2$ outputs $p-1$ if
 $\vecx \in \calV_{\vecg} \cap \overline{\calV_{\vech}}$, $\sigma(\vecx) \ne \vecx$ and $v = \Sigma_{\vecx}$ and 0
 otherwise.


 \paragraph{Grouping Functions.}
      The grouping functions $\phi_U$ and $\phi_V$ are defined as follows (analogous to the so-called ``chessplayer algorithm'' in \cite{r33}):
 \begin{itemize}[label=$\triangleright$,leftmargin=0.5cm]
	\item Grouping $\phi_U$ (corresponding to endpoint in $U$):
\begin{itemize}[leftmargin=0.25cm]
		\item For $\vecx \in \overline{\calV_{\vecg}}$: there exists some $i$ such
		that $\CW_{g_i}(\vecx) = 0$. Consider an edge
		$(\vecx,(a_1,a_2, \dots, a_{m_g}),k)$. We can explicitly list out the multiset containing the monomials
		$t_j = (a_1,a_2, \dots, a_i \gets j, \dots, a_{m_g})$ with multiplicity
	        $t_j(\vecx)$, for each $1 \le j \le L_i$. Since $\CW_{g_i}(\vecx) = 0$, this multiset has size
		multiple of $p$. Hence, we can canonically divide its elements into groups of
		size $p$, counting multiplicities and $\phi_U$ returns the subset containing $(t,k)$.
		\item For $\vecx \in \calV_{\vecg} \cap \overline{\calV_{\vech}}$ such that $\sigma(\vecx) \ne \vecx$: Note that $g_{i}^{p-1}(\vecx) = 0$ for all $i \in \bracks{m_g}$.
		Let $v_1 \in V_1$ be the vertex corresponding to the
		constant monomial $1$. $\phi_U$ groups the edge $(\vecx, v_1, 1)$ (of multiplicity $1$) with the $p-1$ edges $(\vecx, \Sigma_{\vecx}, k)$ for $k \in [p-1]$.
		For any other $t \in  V_1 \setminus \set{v_1}$ and an edge $(\vecx, t,k)$, we have that
		$t = (a_1,\dots, a_{m_g})$ has $a_i \ne 0$ for some $i$. We define
		the multiset containing $t_j = (a_1, \dots, a_{i} \gets j, \dots a_{m_g})$
		with multiplicity $t_j(\vecx)$ for each $1 \le j < L_i$. Since
		$g_{i}^{p-1}(\vecx) = 0$, this multiset has size which is a multiple of $p$, which we can canonically partition into groups of size $p$. Thus, $\phi_U$ on input
		$(\vecx, t, k)$ returns the group containing $(t,k)$.
 	\end{itemize}
	\item Grouping $\phi_V$ (corresponding to endpoint in $V$):
		\begin{itemize}[leftmargin=0.25cm]
		\item For $t \in V_1$ such that $t \neq \prod\limits_{i = 1}^n x_i^{p-1}$:
		there exists a variable $x_i$ with degree less than $p-1$. For
		$\vecx_j = (x_1, \dots, x_{i-1}, x_i \gets j, \dots, x_n)$ with $j \in \bbF_p$
		we define the multiset $\set{(\vecx_j, t(\vecx_j))}_{j \in \bbF_p}$.
		Since $\sum\limits_{j = 0}^{p-1} t(\vecx_j) = 0$, this multiset has
		size multiple of $p$, so we can canonically partition it into groups of size $p$. Then,
		$\phi_V(\vecx, t, k)$ returns the group containing $(\vecx,k)$,
		\item For $v \in V_2$: if $\deg(v) = 0$, then there is no grouping to be done. Else if $\deg(v) \equiv 0 \mod{p}$ then $\phi_V(\vecx,t,k)$ returns
          		$\set{(\vecx,k)}_{\vecx \in v}$.
	  \end{itemize}
\end{itemize}
	Thus, for any vertex with degree $\equiv 0 \mod{p}$, we have provided a grouping function for all its edges. So, for any edge that is not grouped by grouping function at any of its endpoints, then such an endpoint must have degree $\not\equiv 0 \mod{p}$ and hence point to a valid solution of the $\cChevalleyp$ instance.
\end{proof}



\subsubsection{$\cChevalleyp$ is $\PPA_p$--hard}

We show that $\lonely_p \preceq \cChevalleyp$. In the $\cChevalleyp$ instance that we create, we will ensure that there are no solutions of type 0 (as in \autoref{def:combined-chevalley}) and thus, the only valid solutions will be of type 1. In order to do so, we introduce the notions of labeling and proper labeling and prove a generalization of \CWT that we call Labeled \CWT (\autoref{thm:labeledChevalleyWarningTheorem}).

As we will see, the Labeled \CWT, is just a
re-formulation of the original \CWT rather than a generalization. To
understand the Labeled \CWT we start with some
examples that do not seem to satisfy the Chevalley-Warning
condition, but where a solution exists.

\bigskip
\paragr{Example 1.} Consider the case where $p = 3$ and
$f(x_1, x_2) = x_2 - x_1^2$. In this case the Chevalley-Warning condition is not met, since we have $2$ variables and the total degree is also
$2$. But, let us consider a slightly different polynomial where we
replace the variable $x_2$ with the product of two variables
$x_{21}, x_{22}$ then we get the polynomial
$g(x_1, x_{21}, x_{22}) = x_{21} \cdot x_{22} - x_1^2$. Now,
$g$ satisfies \eqref{eq:CWcondition} and hence, we
conclude that the number of roots of $g$ is a multiple of $3$.
Interestingly, from this fact we can argue that there exists a non-trivial solution for $f(\vecx) = 0$.
In particular, the assignment $x_1 = 0$, $x_2 = 0$ corresponds to five
assignments of the variables $x_1$, $x_{21}$, $x_{22}$. Hence, since
$\abs{\calV_{g}} = 0 \mod{3}$, $g$ has another root, which corresponds to a
non-trivial root of $f$. In this example, we applied the
\CWT on a slightly different polynomial than $f$
to argue about the existence of non-trivial solutions of $f$, even though $f$ did not satisfy \eqref{eq:CWcondition} itself.
\bigskip

\paragr{Ignore Some Terms.} The Labeled \CWT formalizes the phenomenon
observed in Example 1 and shows that under certain conditions we can
\textit{ignore some terms} when defining the degree of each polynomial.
For instance, in Example 1, we can ignore the term
$x_1^2$ when computing the degree of $f$ and treat $f$ as a degree $1$
polynomial of $2$ variables, in which case the condition of \CWT is
satisfied.

We describe which terms can be ignored by defining a \textit{labeling}
of the terms of each polynomial in the system. The labels take values in
$\{-1, 0, +1\}$ and our final goal is to ignore the terms with label
$+1$. Of course, it should not be possible to define any labeling that
we want; for example we cannot ignore all the terms of a polynomial.
Next, we describe the rules of a \textit{proper labeling} that will
allow us to prove the Labeled \CWT. We start with
a definition of a labeling.

\begin{definition}[\textsc{Monomial Labeling}]
  \label{def:monomiallabeling}
     Let $\vecf \in \bbF[\vecx]^m$ and let $t_{i j}$ be the $j$-th monomial of the polynomial $f_i \in \bbF[\vecx]$ (written in some canonical sorted order). Let $\calT$ be the
  set of all pairs $(i, j)$ such that $t_{i j}$ is a monomial in
  $\vecf$. A \textit{labeling} of $\vecf$ is a function
  $\lambda : \calT \to \{-1, 0, +1\}$ and we say that $\lambda(i, j)$
  is the \textit{label} of $t_{i j}$ according to $\lambda$.
\end{definition}

\begin{definition}[\textsc{Labeled Degree}] \label{def:labeledDegree}
    For $\vecf \in \bbF[\vecx]^m$ with a labeling $\lambda$,     we define the \textit{labeled degree} of $f_i$ as,
      $ \deg^{\lambda}(f_i)\coloneqq \max_{j\,:\,\lambda(i, j) \ne +1} \deg(t_{i j})$, in words, maximum degree among monomials of $f_i$ labeled either $0$ or $-1$.
\end{definition}

\paragr{Example 1 (continued).} According to the lexicographic ordering,
$f(x_1, x_2) = - x_1^2 + x_2$ and we have the monomials
$t_{1 1} = - x_1^2$ and $t_{1 2} = x_2$. Hence, one possible labeling,
which as we will see later corresponds to the vanilla Chevalley-Warning
Theorem, is $\lambda(1, 1) = \lambda(1, 2) = 0$. According to this labeling,
$\deg^{\lambda}(f) = 2$. Another possible
labeling, that, as we will see, allows us to apply the Labeled
\CWT\!, is $\lambda(1, 1) = +1$ and $\lambda(1, 2) = -1$.
In this case, the labeled degree is $\deg^{\lambda}(f) = 1$.
\bigskip

  As we highlighted before, our goal is to prove the Chevalley-Warning
Theorem, but with the weaker condition that
$\sum_{i = 1}^m \deg^{\lambda}(f_{i}) < n$ instead of
$\sum_{i = 1}^m \deg(f_i) < n$. Of course, we first have to restrict the
space of all possible labelings by defining {\em proper labelings}.
In order to make the condition of proper labelings easier to interpret
we start by defining the notion of a labeling graph.

\begin{definition}[\textsc{Labeling Graph}] \label{def:labelingGraph}
  For $\vecf \in \bbF[\vecx]^m$ with a labeling $\lambda$, we define the
  \textit{labeling graph} $G_{\lambda} = \p{U \cup V, E}$ as a
  directed bipartite graph on vertices
  $U = \set{x_1, \dots, x_n}$ and $V=\set{f_1, \dots, f_m}$. The edge
  $(x_j \to f_i)$ belongs to $E$ if $x_j$ appears in a monomial
  $t_{i r}$ in $f_i$ with label $+1$, i.e. $\lambda(i, r) = +1$. Symmetrically, the edge
  $(f_i \to x_j)$ belongs to $E$ if the $x_j$ appears in a monomial
  $t_{i r}$ in $f_i$ with label $-1$, i.e. $\lambda(i, r) = -1$.
\end{definition}

\paragr{Example 2.} Let $p = 2$ and
$f_1(x_1, x_2, x_3, x_4) = x_1 x_2 - x_3$,
$f_2(x_1, x_2, x_3, x_4) = x_1 x_3 - x_4$. In this system, if we use the
lexicographic monomial ordering we have the monomials
$t_{1 1} = x_1 x_2$,
$t_{1 2} = - x_3$, $t_{2 1} = x_1 x_3$, $t_{2 2} = - x_4$. The following
figure shows the graph $G_{\lambda}$ for the labeling
$\lambda(1, 1) = +1$, $\lambda(1, 2) = -1$, $\lambda(2, 1) = +1$ and
$\lambda(2, 2) = -1$.

\begin{center}
\begin{tikzpicture}
	\tikzset{
		inner sep=0,outer sep=3,
		neg/.style={color=Gred},
		posi/.style={color=Gblue}}
	\def\hgap{1.5}
	\def\vgap{1.5}
	\node (f1) at (0,\vgap) {$f_1$};
	\node (f2) at (\hgap,\vgap) {$f_2$};

	\node (x1) at (-\hgap,0) {$x_1$};
	\node (x2)at (0,0) {$x_2$};
	\node (x3) at (\hgap,0) {$x_3$};
	\node (x4) at (2*\hgap,0) {$x_4$};

	\path[-{Stealth[length=5pt]},line width=.4pt,black]
	(f1) edge[neg] (x3)
	(f2) edge[neg] (x4)
	(x1) edge[posi] (f1)
	(x2) edge[posi] (f1)
	(x1) edge[posi] (f2)
	(x3) edge[posi] (f2);
\end{tikzpicture}
\end{center}


\begin{definition}[\textsc{Proper Labeling}] \label{def:properLabel}
  Let $\vecf \in \bbF[\vecx]^m$ with a labeling $\lambda$.
    We say that the labeling $\lambda$ is \textit{proper} if the following
  conditions hold.\\[-6mm]
  \begin{enumerate}
    \item[(1)] For all $i$, either $\lambda(i, j) \in \set{-1, 1}$ for all $j$, or $\lambda(i, j) = 0$ for all $j$.
    \item[(2)] If two monomials $t_{i j}$, $t_{i j'}$ contain the same
          variable $x_k$, then $\lambda(i, j) = \lambda(i, j')$.
    \item[(3)] If $\lambda(i, j) = -1$, then $t_{i j}$ is multilinear.
    \item[(4)] If $x_k$ is a variable in the monomials
          $t_{i j}$, $t_{i' j'}$,
          with $i \neq i'$ and $\lambda(i, j) = -1$, then
          $\lambda(i', j') = +1$.
    \item[(5)] If $\lambda(i, j) \neq 0$, then there exists a $j'$ such that
          $\lambda(i, j') = -1$.
    \item[(6)] The labeling graph $G_{\lambda}$ contains no directed cycles.
  \end{enumerate}
\end{definition}

\noindent We give an equivalent way to understand the definition of a proper labeling.
\begin{itemize}[label=$\triangleright$]
\item Condition (1) : there is a partition of the polynomial system $\vecf$ into polynomial systems $\vecg$ and $\vech$ such that all monomials in $\vecg$ are labeled in $\set{+1, -1}$ and all monomials in $\vech$ are labeled $0$.
\item Condition (2) : each polynomial $g_i$ in $\vecg$ can be written as $g_i = g_i^+ + g_i^{-}$, such that $g_i^+$ and $g_i^-$ are polynomials on a disjoint set of variables.
\item Condition (3) : Each $g_i^-$ is multilinear.
\item Condition (4) : Any variable $x_k$ can appear in at most one of the $g_i^-$. Moreover, if an $x_k$ appears in some $g_i^-$, it does not appear in any $h_j$ in $\vech$.
\item Condition (5) : Every $g_i^-$ involves at least one variable.
\item Condition (6) : The graph $G_{\lambda}$ is essentially between polynomials in $\vecg$ and the variables that appear in them, with an edge $(g_i \to x_k)$ if $x_k$ appears in $g_i^-$ or an edge $(x_k \to g_i)$ if $x_k$ appears in $g_i^+$.
\item Note that $\deg^\lambda(g_i) = \deg(g_i^-)$, whereas $\deg^\lambda(h_j) = \deg(h_j)$.
\end{itemize}

\noindent It is easy to see that the trivial labeling $\lambda(i, j) = 0$ is always proper.
As we will see this special case of the Labeled \CWT\! corresponds to the
original \CWT\!. Note that in this case the
labeling graph $G_{\lambda}$ is an empty graph. Also, given a
system of polynomials $\vecf$ and a labeling $\lambda$, it is possible to
check in polynomial time whether the labeling $\lambda$ is proper or not.

\bigskip
\paragr{Example 2 (continued).} It is an instructive exercise to verify that the labeling $\lambda$ specified was indeed a proper labeling of $\vecf$.
\bigskip


\begin{theorem}[\textsc{Labeled Chevalley-Warning Theorem}]
  \label{thm:labeledChevalleyWarningTheorem}
    Let $\F_q$ be a finite field with characteristic $p$ and
  $\vecf \in \F_q[\vecx]^m$. If $\lambda$ is a proper labeling of
  $\vecf$ with $\sum_{i = 1}^m \deg^{\lambda}(f_i) < n$, then $|\calM_{\vecf}| = 0$. In particular, $\abs{\calV_{\vecf}} \equiv 0 \mod{p}$.
\end{theorem}
\begin{proof}
We can re-write $\CW_{\vecf}(x)$ as
$\sum_{S \subseteq [m]} \prod_{i \in S} (-1)^{|S|} f_i^{p-1}$.
We'll show that every monomial appearing in the expansion of
$\prod_{i \in S} f_i^{p-1}$ will have at least one variable with degree at most
$p-1$. For simplicity, we focus on the case $S = [m]$ and the other cases of
$S$ follow similarly.

We index a monomial of $\prod_{i \in [m]} f_i^{p-1}$ with a tuple
$((j_{11},j_{12}, \dots, j_{1(p-1)}), \dots, (j_{m1}, \dots, j_{m(p-1)}))$ with
$1 \leq j_{i\ell} \leq L_{i}$ where $L_i$ is the number of monomials in the
explicit representation of $f_i$.
The coordinates $(j_{i1}, \dots, j_{i(p-1)})$ represent the indices of the
monomials chosen from each of the $p-1$ copies of $f_i^{p-1}$. More succinctly,
we have $t = \prod_{i=1}^m \prod_{\ell=1}^{p-1} t_{i,j_{i\ell}}$.\\[-0.8mm]

\noindent {\em Case 1. $\lambda(i,j_{i\ell}) \in \set{0, -1}$, for all $(i,\ell)$:}\\
Here, $\deg(t) \le (p-1) \sum_{i=1}^m \deg^{\lambda}(f_i)$ which, by our assumption, is strictly less than $(p-1)n$. Hence, there is a variable with degree less than $p-1$.\\[-0.8mm]

\noindent {\em Case 2. There is a unique $i$ with $\lambda(i,j_{i\ell}) = +1$ for some $\ell$:} (warmup for case 3)\\
That is, for all $i'\neq i$, $\lambda(i',j_{i'\ell})  \in \set{0, -1}$. By condition (5) of proper labeling there exists a $j' \neq j_{i\ell}$ such that  $\lambda(i,j') = -1$. Let $x_k$ be a variable in the monomial $t_{ij'}$. By condition (2), $x_k$ is not present in the monomial $t_{i,j_{i\ell}}$ and by condition (3), its degree in $(t_{i,j_{i,1}}, \dots, t_{i,j_{i,p-1}})$ is at most $p-2$. Additionally, by condition (4), any monomial of $f_{i'}$ for $i' \neq i$ containing $x_k$ must have label $+1$, but $\lambda(i',j_{i',\ell})$ are all in $\set{0, -1}$. Hence, $x_k$ does not appear in any other monomial of $t$ and its degree on $t$ is equal to its degree in $(t_{i,j_{i,1}} \cdots t_{i,j_{i,p-1}})$, which is strictly less than $p-1$.\\[-0.8mm]

\noindent {\em Case 3. $I = \set{i : \lambda(i,j_{i\ell}) = +1 \text{ for some } \ell}$:}\\
In the labeling graph $G_\lambda$, let $i \in I$ be such that there is no path from $f_i$ to any other $f_{i'}$ for $i' \in I$. Such an $i$ exists due to acyclicity of $G_{\lambda}$, i.e. condition (6). Let $\ell$ be such that $\lambda(i,j_{i\ell}) = +1$. Again, by condition (5) of proper labeling there exists a $j' \neq j_{i\ell}$ such that  $\lambda(i,j') = -1$. Let $x_k$ be a variable in the monomial $t_{ij'}$. By condition (2), $x_k$ is not present in the monomial $t_{i,j_{i\ell}}$ and by condition (3), its degree in $(t_{i,j_{i,1}}, \dots, t_{i,j_{i,p-1}})$ is at most $p-2$. Additionally, by condition (4), any monomial of $f_{i'}$ for $i' \neq i$ containing $x_k$ must have label $+1$. For $i' \notin I$, $\lambda(i',j_{i',\ell})$ are all in $\set{0, -1}$. And for $i' \in I$, variable $x_k$ cannot appear with $+1$ label in $f_{i'}$ by our choice of $f_i$. Hence, $x_k$ does not appear in any other monomial of $t$ and its degree on $t$ is equal to its degree on $(t_{i,j_{i,1}} \cdots t_{i,j_{i,p-1}})$, which is strictly less than $p-1$.

\end{proof}


\medskip
\noindent We are now ready to prove the $\PPAp$-hardness of $\cChevalleyp$.

\begin{lemma} \label{lem:combinedChevalleyPPApHardness}
    For all primes $p$, $\cChevalleyp$ is $\PPAp$-hard.
\end{lemma}

\begin{proof}    We prove that $\lonelyp \reducible \cChevalleyp$. Let us assume (without loss of
    generality from \autoref{lem:restricted-lonely}) that the $\lonelyp$
    instance has a single distinguished vertex represented by $0^n$. We'll assume that $0^n$ is isolated, otherwise, no further reduction is necessary.\medskip

    \paragr{Pre-processing.} We slightly
    modify the given circuit $\calC$ by defining
    $\calC': \bbF_p^n \to \bbF_p^n$ as follows:
    \[ \calC'(v) = \left\{
                \begin{array}{cl}
                  v         &     \text{, if } \calC^p(v) \neq v \\
                  \calC(v)  &     \text{, otherwise}
                \end{array}
                \right.
    \]
    Since $p$ is a prime, a vertex $v \in \bbF_p^n$ has $\deg(v) = 1$ if and only if
    $\calC^p(v) = v$ and $\calC(v) \neq v$. By modifying the circuit, we changed this
    condition to just $\calC'(v) \neq v$, which facilitates our reduction.

    Circuit
    $\calC'$ is composed of the $\bbF_p$-addition $(+)$, $\bbF_p$-multiplication $(\times)$
    and the constant $(1)$ gates. However, we require the input of $\cChevalleyp$ to be a
    zecote polynomial system, and so we further modify the circuit $\calC'$ to eliminate all the constant $(1)$ gates, without changing it's behavior -- this is possible because we assume $\calC'(0^n) = 0^n$.
                \begin{claim}
        \label{cl:circuitWithout1}
          Given circuit $\calC'$ with $(+,\times, 1)$ gates, there exists circuit
          $\bar{\calC}$ with $(+, \times)$ gates such that
          \[ \bar{\calC}(\vecv) = \infork{0^n & \text{, if } \vecv = 0^n \\ \calC'(\vecv) & \text{, otherwise}}
          \]
    \end{claim}
    \begin{proof}[Proof of \autoref{cl:circuitWithout1}]
    	We replace all instances of the $(1)$ gate by the function $\mathds{1}_{\set{\vecv \ne 0^n}}$, which we can compute using only $(+, \times)$ gates as follows: For any $x, y \in \bbF_p$, observe that $\mathds{1}_{\set{x \ne 0}} \vee \mathds{1}_{\set{y \ne 0}} = x^{p-1} + y^{p-1} - x^{p-1} y^{p-1}$. We can thus recursively compute $\bigvee_{i=1}^n \mathds{1}_{\set{v_i \ne 0}}$ using only $(+, \times)$ gates. Thus, $\bar{\calC}(\vecv) = \calC'(\vecv)$ for all $\vecv \ne 0^n$. And $\bar{\calC}(0^n) = 0^n$, since $\bar{\calC}$ is computed with only $(+, \times)$ gates.
    \end{proof}

        Thus, we can transform our original circuit $\calC$ into a circuit $\bar{\calC}$ with just $(+,\times)$ gates. For simplicity, we'll write $\bar{\calC}$ as simply $\calC$ from now on.\medskip

        As an intermiate step in the reduction we describe a system of
        polynomials $\vecf_{{\calC}}$ over $2n+s$ variables
        $(x_1, \dots, x_n, z_1, \dots, z_s, y_{1}, \dots, y_{n})$, where $s$ is the size
        of the circuit $\calC$. The variables
        $\vecx = (x_1, \dots,x_n)$ correspond to the input of $\calC$, the
        variables $\vecy = (y_1, \dots, y_n)$ correspond to the output and the variables
        $\vecz = (z_1, \dots, z_s)$  correspond to the gates of $\calC$. For an addition
        gate $(+)$ we include a polynomial of the form
        \[           f(a_1, a_2, a_3) =  a_2 + a_3 - a_1, \]
        where $a_1$ is the variable corresponding to the output of the $(+)$ gate and $a_2,a_3$
        are the variables corresponding to its two inputs. Similarly for a multiplication
        $(\times)$ gate, we include a polynomial of the form
        \[           f(a_1,a_2,a_3) = a_2 \cdot a_3 -  a_1 \]
        Finally, for the output of the circuit, we include the polynomial
        \[              f(a,y_i) =  a -  y_i, \]
        where $a$ is the variable corresponding to the $i$-th output gate of
        $\calC$.
        It holds that
        \[  \calC(\vecx) = \vecy \quad \Longleftrightarrow \quad
         \vecf_{\calC}(\vecx,\vecy,\vecz) = \boldsymbol{0}.\]

		\noindent We now describe the reduction from an instance of $\lonely_p$ to that of $\cChevalley_p$. In order to do this, we need to specify a system of polynomials $(\vecg, \vech)$ and a permutation $\sigma$ such that $|\sigma| = p$. In addition, we will provide a proper labeling $\lambda$ for $\vecg$ satisfying the degree condition. We will also ensure that $\a{\sigma}$ acts freely on $\calV_{\vecg} \cap \overline{\calV_{\vech}}$. And hence, the only valid solutions for the resulting $\cChevalley_p$ instance will be $\vecx \in \calV_{\vecg} \cap \calV_{\vech}$.

        \paragraph{Definition of $\vecg$.} The polynomial system $\vecg$ contains the
              following systems of polynomials.
             \[
              \begin{array}{c}
                   \vecf_{\calC}(\vecx_1, \vecx_2, \vecz_{1,2})\\
                   \vecx_2 - \vecx_3 \\
                   \vecf_{\calC}(\vecx_3, \vecx_4, \vecz_{3,4})\\
                   \vecx_4 - \vecx_5 \\
                   \vdots \\
                   \vecf_{\calC}(\vecx_{2p-1}, \vecx_{2p}, \vecz_{2p-1,2p})
              \end{array}
              \]
              Note that there are $N = (2n+s)p$ variables in total.


        \medskip
        \paragraph{Labeling $\lambda$ of $\vecg$.} For the polynomials belonging to a
        system  of the form $\vecf_{\calC}$, the labeling is equal to
        $-1$ for the monomials corresponding to the output of each gate
        and $+1$, otherwise. For instance, let $a_2 + a_3 - a_1$ be the $i$-th
        polynomial of $\vecg$ corresponding to a $(+)$ gate and let
        $a_1 \prec a_2 \prec a_3$, then
        $\lambda(i,1) = -1$ and $\lambda(i,2) = \lambda(i,3) = +1$.

        For the polynomials belonging to a system of the form $\vecx_i - \vecx_{i+1}$,
        the labeling is equal to $-1$ for the monomials with variables in $\vecx_{i+1}$
        and $+1$ for the monomials with variables in  $\vecx_{i}$.

        \begin{claim}
          \label{cl:properLabelingForf1}
          The labeling $\lambda$ for $\vecg$ is proper.
        \end{claim}

\begin{proof}[Proof of \autoref{cl:properLabelingForf1}]
        By \autoref{def:properLabel}, the labeling $\lambda$ is proper if the following conditions hold.
      \begin{description}
        \item[Condition 1.] \emph{For all $i$, either $\lambda(i, j) \in \set{-1, 1}$ for all $j$, or $\lambda(i, j) = 0$ for all $j$.} \\
              In the labeling $\lambda$, there are no labels equal to
              $0$, so this condition holds trivially.
        \item[Condition 2.] \emph{If two monomials $t_{i j}$, $t_{i j'}$ contain the same
              variable $x_k$, then $\lambda(i, j) = \lambda(i, j')$.}\\
              By construction of $\vecg$, no variable appears twice in the same
              polynomial with a different labeling. For polynomials of
              $\vecf_{\calC}$, this holds because the output variable of a gate is
              not simultaneously an input variable and all
              input variables have the same labeling. For polynomials in a system of
              the form $\vecx_i - \vecx_{i+1}$, each polynomial contains two different
              variables.
        \item[Condition 3.] \emph{If $\lambda(i, j) = -1$, then $t_{i j}$ is multilinear.} \\
              For polynomials of $\vecf_{\calC}$, only the output variable of a
              gate has label $-1$ and by definition this monomial is linear. For
              polynomials in a system of the form $\vecx_i - \vecx_{i+1}$, all monomials
              are linear, so the condition holds trivially.
        \item[Condition 4.] \emph{If $x_k$ is a variable in the monomials $t_{i j}$,
              $t_{i' j'}$,
              with $i \neq i'$ and $\lambda(i, j) = -1$, then $\lambda(i', j') = +1$.} \\
              Observe that all monomials with label $-1$ contain only a single
              variable, so we refer to a monomial $x_k$ with label $-1$.
              For a polynomial in $\vecf_{\calC}$, a monomial $x_k$ with label $-1$
              corresponds to the output of a gate. Hence, if $x_k$ appears in other
              monomials of $\vecf_{\calC}$, these monomials correspond to inputs
              and have label $+1$. Also, if $x_k$ is an output variable of
              $\vecf_{\calC}$, then it might appear in a polynomial of the form
              $a_1 - a_2$. However, by construction the monomials of $x_i - x_{i+1}$ that
              correspond to output variables of $\vecf_{\calC}$ have label $+1$.
        \item[Condition 5.] \emph{If $\lambda(i, j) \neq 0$, then there exists a $j'$ such
              that $\lambda(i, j') = -1$.} \\
              By the definition of $\lambda$, all polynomials of $\vecg$ have a monomial
              with label $-1$. These are the monomials that correspond to the outputs of
              a gate for the systems of the form $\vecf_{\calC}$ and the monomials
              that correspond to $\vecx_{i+1}$ for the systems of the form
              $\vecx_{i}- \vecx_{i+1}$.
        \item[Condition 6.] \emph{The labeling graph $G_{\lambda}$ contains no cycles.}\\
              Each system of the form $\vecx_i - \vecx_{i+1}$ has incoming
              edges with variables appearing only in the $i$-th copy of
              $\vecf_{\calC}$ and outgoing edges with variables
              appearing only in the $(i+1)$-th copy of $\vecf_{\calC}$. Also, the
              variables appearing on the $i$-th copy of $\vecf_{\calC}$ might
              appear only in the systems $\vecx_{i-1} - \vecx_i$ and
              $\vecx_i - \vecx_{i+1}$. Hence, $G_{\lambda}$ has no cycles that contain
              vertices of two different copies of $\vecf_{\calC}$ or of a copy of
              $\vecf_{\calC}$ and a system of the form $\vecx_{i-1} - \vecx_i$.

              It is left to argue that the labeling graph restricted to a copy of
              $\vecf_{\calC}$  does not have any cycles. Let the vertices of
              $\vecf_{\calC}$ be ordered according to the topological ordering of
              $\calC$. This restricted part of $G_{\lambda}$ corresponds exactly to
              the graph of $\calC$, which by definition is a DAG. Hence, $G_{\lambda}$
              contains no cycles.\qedhere
      \end{description}
\end{proof}

        We also need to show that for this labeling $\vecg$ satisfies the labeled
        Chevalley condition.
        \begin{claim}
          \label{cl:ChevalleyForf1}
          The labeled Chevalley condition
          $\sum_{i = 1}^{m_g} \deg^{\lambda}(g_i) < N$ holds for $\vecg$ with
          labeling $\lambda$.
        \end{claim}
        \begin{proof}              Each polynomial of $\vecg$ has a unique monomial with $\lambda(i,j) = -1$
              and this monomial has degree $1$. Thus,
              $\sum_{i = 1}^{m_g} \deg^{\lambda}(g_{i}) = m_g$.
              On the other hand, the $i$-th polynomial of $\vecg$ has exactly one
              variable that has not appeared in any of the previous polynomials. More
              specifically, the number of variables is equal to $m_g + n$, where $n$
              is the size of the input of $\calC$. Hence, the labeled Chevalley
              condition holds for $\vecg$.
        \end{proof}

        \paragraph{Definition of $\vech$.} The system of polynomials $\vecg$ allows
        us to compute the $p$ vertices given by ${\calC}^{i}(\vecx)$ for
        $i \in \bracks{p+1}$. From the definition of $\lonelyp$ and our pre-processing on
        $\calC$, this group of $p$ vertices is a hyperedge if and only if
        $\calC(\vecx) \neq \vecx$. Since solutions of $\lonelyp$ are lonely
        vertices, we define $\vech$ to exclude $\vecx$ such that
        $\calC(\vecx) \neq \vecx$. Namely, we set $\vech$ to be the system of
        polynomials
                          \[ \vecx_1 - \vecx_2.\]

        \paragraph{Definition of permutation $\boldsymbol{\sigma}$.} In the description
        of $\vecf = (\vecg, \vech)$, we have used the following vector of variables:
                    \[  \vecx = (\vecx_1, \vecx_2, \dots, \vecx_{2p},
                          \vecz_{1,2}, \vecz_{3,4}, \dots, \vecz_{2p-1,2p} )
                    \]
      	We define the permutation $\sigma$ such that $\sigma(\vecx) = (\vecx_3, \vecx_4, \dots, \vecx_{2p},
                \vecx_1, \vecx_2, \vecz_{3,4}, \vecz_{5,6}, \dots, \vecz_{2p-1,2p},
                \vecz_{1,2})$, as illustrated in the following figure. The \textcolor{Gblue}{blue} arrows indicate the polynomials $\vecg$ and the \textcolor{Ggreen}{green} arrows indicate the permutation $\sigma$ in the case of $p=3$.

		\begin{center}
		\begin{tikzpicture}
			\tikzset{
				zij/.style = {inner sep=0pt, outer sep=2pt, midway, above=0.1cm},
				comp/.style = {{Stealth[length=5pt]}-, draw=Gblue, thick}
			}
			\def\hgap{4.5}
			\def\vgap{2.3}
			\node (x1) at (0,0) {$\vecx_1$};
			\node (x2) at (\hgap, 0) {$\vecx_2$} edge[comp] node[zij] (z12) {$\vecz_{1,2}$} (x1);
			\node (x3) at (0,-\vgap) {$\vecx_3$} edge[comp] node[midway, sloped, above] (eq23) {$=$} (x2);
			\node (x4) at (\hgap, -\vgap) {$\vecx_4$} edge[comp] node[zij] (z34) {$\vecz_{3,4}$} (x3);
			\node (x5) at (0,-2*\vgap) {$\vecx_5$} edge[comp] node[midway, sloped, above] (eq45) {$=$} (x4);
			\node (x6) at (\hgap, -2*\vgap) {$\vecx_6$} edge[comp] node[zij] (z56) {$\vecz_{5,6}$} (x5);

			\path[-{Stealth[length=5pt]},line width=0.7pt,Ggreen]
			(x3) edge (x1)
			(x5) edge (x3)
			(x1) edge[out=-120,in=120] (x5)
			(x4) edge (x2)
			(x6) edge (x4)
			(x2) edge[out=-60,in=60] (x6)
			([shift={(0,0.1)}]z34.east) edge[out=60,in=-60] ([shift={(0,-0.1)}]z12.east)
			([shift={(0,0.1)}]z56.east) edge[out=60,in=-60] ([shift={(0,-0.1)}]z34.east)
			([shift={(0,-0.1)}]z12.west) edge[out=-120,in=120] ([shift={(0,0.1)}]z56.west);
		\end{tikzpicture}
		\end{center}

          \begin{claim}
            \label{cl:sigmaActsFreelyAndOrderp}
            The group $\a{\sigma}$ has order $p$ and acts freely on
            $\calV_{\vecg} \cap \overline{\calV_{\vech}}$.
          \end{claim}

          \begin{proof}              In order to see that $\abs{\sigma} = p$, note that the input of $\sigma$ consists
              of $3p$ blocks of variables. The permutation $\sigma$ performs a rotation of the
              first $2p$ blocks by two positions and of the last $p$ blocks by one
              position.
              All that remains is to show that $\a{\sigma}$ acts freely on
                    $\calV_{\vecg} \cap \overline{\calV_{\vech}}$.
                    First, we show that $\a{\sigma}$ defines a group action on
                    $\calV_{\vecg} \cap \overline{\calV_{\vech}} $, that is
                    for all $\vecx \in \calV_{\vecg} \cap
                    \overline{\calV_{\vech}}$, it holds that $\sigma(\vecx) \in
                    \calV_{\vecg} \cap \overline{\calV_{\vech}}$.
                    Let $\vecx = (\vecx_1, \vecx_2,
                    \dots, \vecx_{2p-1}, \vecx_{2p}, \vecz_{1,2}, \vecz_{3,4}, \dots, \vecz_{2p-
                    1,2p} ) \in \calV_{\vecg} \cap \overline{\calV_{\vech}} $,
                    then
                    \begin{itemize}
                        \item $\vecx \in \calV_{\vecg}$ implies that
                        $\vecf_{\calC}(\vecx_{2i-1},\vecx_{2i},
                        \vecz_{2i-1,2i}) = 0$ for $i \in \bracks{p}$ and
                        $\vecx_{2i} = \vecx_{2i+1}$ for $i \in \bracks{p-1}$
                        \item $\vecx \in \overline{\calV_{\vech}}$ implies that
                        $\vecx_1 \neq \vecx_2$, that is,
                        $\calC(\vecx_{1}) \neq \vecx_1$ since
                        $\vecf_{\calC}(\vecx_{1},\vecx_{2},
                        \vecz_{1,2}) = 0 \Leftrightarrow \vecx_2 =
                        \calC(\vecx_{1}) $.
                    \end{itemize}
                Now, $\sigma(\vecx) = (\vecx_{3}, \vecx_{4},
                \ldots, \vecx_{1}, \vecx_{2}, \vecz_{3,4}, \vecz_{5,6}, \ldots,
                \vecz_{1,2}) \in  \calV_{\vecg} \cap
                \overline{\calV_{\vech}}$ holds because
                \begin{itemize}
                    \item $\vecf_{\calC}(\vecx_{2i-1},\vecx_{2i},
                    \vecz_{2i-1,2i}) = 0$ for $i \in \bracks{p}$ and
                    $\vecx_{2i} = \vecx_{2i+1}$ for $i \in \bracks{p-1}$, which holds
                    from $\vecx \in \calV_{\vecg}$. Additionally, $\vecx_1 = \vecx_{2p}$ holds because we pre-processed $\calC$ such that $\calC^p(\vecx_1) = \vecx_{1}$,
                    \item $\vecx_{3} \neq \vecx_{4}$, which holds because
                    $\vecx_{4}  = \calC(\vecx_{3})$ for $i \in \bracks{p}$
                    and from the definition of $\calC$, $\calC(\vecx_{1})
                    \neq \vecx_1$ implies that $\vecx_{2i} \neq \vecx_{2i-1}$ for all
                    $i \in \bracks{p}$.
                \end{itemize}

              Finally, if $\vecx \in \calV_{\vecg} \cap
              \overline{\calV_{\vech}}$, by
              construction of $\cal{C}$, we have that $\vecx_{2k} \neq \vecx_{2j}$ for
              $k\neq j$ and thus $\sigma(\vecx) \ne \vecx$ simply because $\vecx_3 \ne \vecx_1$. Thus, we conclude that $\a{\sigma}$ acts freely on
              $\calV_{\vecg} \cap \overline{\calV_{\vech}}$.
          \end{proof}

        		\paragraph{Putting it all together.} The solution of this instance of $\cChevalleyp$ cannot be a vector
                $\vecx \in \calV_{\vecg} \cap \overline{\calV_{\vech}}$  with
                $\sigma(\vecx)
                \not\in \calV_{\vecg} \cap \overline{\calV_{\vech}}$ or
                $\sigma(\vecx) = \vecx$, since we know from
                \autoref{cl:sigmaActsFreelyAndOrderp} that $\a{\sigma}$ acts freely on
                $\calV_{\vecg} \cap \overline{\calV_{\vech}}$.
                We also have from \autoref{thm:labeledChevalleyWarningTheorem} that the solution also cannot be a max-degree monomial in the expansion of $\CW_{\vecg}(\bx) = \prod (1-g_i^{p-1})$.
                Thus, the solution must be an $\vecx \neq \vec{0}$ such that
                $\vecf(\vecx) = \vec{0}$. Let $\vecx_1$ denote the first $n$
                coordinates of $\vecx$, then $\vecf(\vecx) = \vec{0}$ implies that
                $\vecx_1 = \calC(\vecx_1)$ and $\vecx \neq \vec{0}$ implies that
                $\vecx_1 \neq \vec{0}$. Hence, $\vecx_1$ corresponds to
                a lonely vertex of the $\lonelyp$ instance.
\end{proof} 
